{
This paper lays the foundation for building a general-purpose task-aware retriever that can follow natural language instructions. 
We introduced a new problem, \emph{retrieval with instructions}, to model users' intents explicitly.
We presented \rib, the first large-scale retrieval dataset with expert-written annotations. 
Building upon \rib, we trained the first instruction-following retrieval system by massive multi-task instruction-tuning, \model, adopting two widely used architectures. 
\model advances the state of the art on the popular zero-shot retrieval benchmarks BEIR and LOTTE as well as on our newly introduced challenging evaluation setup, \task. 
Our analysis shows that key factors to building a successful multi-task instruction-following retrieval system include informative instructions at training and test time, diversity in data and model scale, and carefully designed negative samples. 
We conclude with two interesting open questions, which future work can explore.  

{
\paragraph{Improving efficiency of instruction-following retrievers. } Although our \model-full model shows the effectiveness of instruction-tuning for retrieval, \model-dual shows some performance drop from its non-instruction-following counterpart, caused by large performance deteriorations on a few datasets. 
We hypothesize that a smaller model size (i.e., 110 million parameters) and limited interactions between query and document embeddings make building an instruction-following retrieval system more challenging. 
Future work can study the effect of scaling up bi-encoder models  as well as explore modeling architectures that enable rich interaction but are still more efficient than the cross-encoder, such as ColBERT-v2~\cite{santhanam-etal-2022-colbertv2}. 
}
}


\paragraph{Further scaling up the number of datasets.}
Retrieval tasks are excluded in prior work on instruction-following of LLMs. 
This work is the first to explore instruction-tuning in the area of retrieval, and we annotate more than 100 instructions for approximately 40 tasks, and we demonstrate the effectiveness of the dataset scale in retrieval. 
Recent work~\cite{wang2022benchmarking,flan_palm} show that scaling up the number of the training datasets improves LLMs' ability to adapt to new task via instructions. 
We open-source our instruction data and call for community efforts to collect more retrieval tasks and human-written instructions as in instruction-following for LMs~\cite{wang2022benchmarking,bach2022promptsource}, to investigate whether further increasing the number of the datasets (e.g., more than 100 datasets) improves zero-shot and cross-task retrieval. 



